\subsection{Simpson's Paradox}
\totalscore{25}
You are investigating the effectiveness of an experimental drug against a new deadly disease.
You are given access to data collected by the RIVM (National Institute for Public Health and the Environment). 
You divide the diseased patients into two groups: those that took the drug, and those that didn't take the drug.
Some of the patients recovered, others unfortunately didn't recover and died. 
The reasons why some patients were treated and others were not, are unknown to you.
You find the following numbers:\footnote{Luckily, these numbers are made up.}
\begin{center}\begin{tabular}{l|rr|r|r}
          & Recovery & No recovery & Total & Recovery rate \\
  \hline
  Drug    & 2000       & 2000          & 4000    & 50\%       \\
  No drug & 1600       & 2400          & 4000    & 40\%       \\
  \hline
  Total   & 3600       & 4400          & 8000    &               \\
\end{tabular}\end{center}
Upon closer inspection of the data, you notice something peculiar when you group patients according to sex.
\begin{center}\begin{tabular}{l|rr|r|r}
  Males & Recovery & No recovery & Total & Recovery rate \\
  \hline
  Drug           & 1800       & 1200          & 3000    & 60\%       \\
  No drug        & 700        & 300           & 1000    & 70\%       \\
  \hline
  Total          & 2500       & 1500          & 4000    &               \\
\end{tabular}

\medskip
\begin{tabular}{l|rr|r|r}
  Females & Recovery & No recovery & Total & Recovery rate \\
  \hline
  Drug             & 200        & 800           & 1000    & 20\%       \\
  No drug          & 900        & 2100          & 3000    & 30\%       \\
  \hline
  Total            & 1100       & 2900          & 4000    &               \\
\end{tabular}\end{center}

Interestingly, the probability that a male patient recovers is lower for the treated male patients than for the untreated male patients, and we observe a similar pattern for the female patients.
Yet, if we look at all patients (not conditioning on sex), then the probability to recover is higher for the treated patients than for the untreated.
This phenomenon is an example of ``Simpson's paradox''. It stems from our tendency to interpret the observed correlations as causation.

\begin{figure}\centering
\begin{tikzpicture}
  \begin{scope}
    \node[ndout] (X) at (0,0) {$X$};
    \node[ndout] (Y) at (2,0) {$Y$};
    \node[ndout] (Z) at (1,1) {$Z$};
    \draw[tuh] (X) to (Y);
    \draw[tuh] (Z) to (X);
    \draw[tuh] (Z) to (Y);
    \node at (1,-0.75) {(a)};
  \end{scope}
%  \node at (11,1) {\begin{tabular}{ll}
%  $X$ & treatment \\
%  $Z$ & sex \\
%  $Y$ & outcome \\
%  $L$ & unobserved confounder
%  \end{tabular}};
  \begin{scope}[xshift=4cm]
    \node[ndout] (X) at (0,0) {$X$};
    \node[ndout] (Y) at (2,0) {$Y$};
    \node[ndout] (Z) at (1,1) {$Z$};
%    \node[ndlat] (L) at (0,2) {$L$};
    \draw[tuh] (X) to (Y);
%    \draw[tuh] (L) to (Z);
%    \draw[tuh] (L) to (X);
    \draw[huh] (X) edge (Z);
    \draw[tuh] (Z) to (Y);
    \node at (1,-0.75) {(b)};
  \end{scope}
  \begin{scope}[xshift=0cm,yshift=-3cm]
    \node[ndout] (X) at (0,0) {$X$};
    \node[ndout] (Y) at (2,0) {$Y$};
    \node[ndout] (Z) at (1,1) {$Z$};
%    \node[ndlat] (L1) at (0,2) {$L_1$};
%    \node[ndlat] (L2) at (2,2) {$L_2$};
    \draw[tuh] (X) to (Y);
%    \draw[tuh] (L1) to (Z);
%    \draw[tuh] (L1) to (X);
    \draw[huh] (X) edge (Z);
%    \draw[tuh] (L2) to (Z);
%    \draw[tuh] (L2) to (Y);
    \draw[huh] (Z) edge (Y);
    \node at (1,-0.75) {(c)};
  \end{scope}
  \begin{scope}[xshift=4cm,yshift=-3cm]
    \node[ndout] (X) at (0,0) {$X$};
    \node[ndout] (Y) at (2,0) {$Y$};
    \node[ndout] (Z) at (1,1) {$Z$};
    \draw[tuh] (X) to (Y);
    \draw[tuh] (X) to (Z);
    \draw[tuh] (Z) to (Y);
    \node at (1,-0.75) {(d)};
  \end{scope}
  \begin{scope}[xshift=8cm,yshift=-3cm]
%      \draw[shadow] (-1,-2) rectangle (3,2);
    \node[ndout] (X) at (0,0) {$X$};
    \node[ndout] (Y) at (2,0) {$Y$};
    \node[ndout] (Z) at (1,1) {$Z$};
%    \node[ndlat] (L) at (1,-1) {$L$};
    \draw[tuh] (X) to (Y);
    \draw[tuh] (Z) to (X);
    \draw[tuh] (Z) to (Y);
%    \draw[tuh] (L) to (X);
%    \draw[tuh] (L) to (Y);
    \draw[huh,bend right] (X) edge (Y);
    \node at (1,-0.75) {(e)};
  \end{scope}
\end{tikzpicture}
\caption{\label{fig:simpson_paradox_causal_hypotheses}Different causal hypotheses (observable ADMGs of an L-CBN or SCM) to explain the patient data.}
\end{figure}

In this exercise we will make use of causal reasoning with L-CBNs or SCMs to resolve Simpson's paradox.
Consider the following binary variables:
\begin{center}\begin{tabular}{llll}
  symbol & meaning & domain & meaning \\
  \hline
  $X$ & treatment & $\C{X} = \{0,1\}$ & no drug/drug \\
  $Y$ & outcome   & $\C{Y} = \{0,1\}$ & died/recovered \\
  $Z$ & sex       & $\C{Z} = \{0,1\}$ & male/female
\end{tabular}\end{center}
We will consider different possible causal hypotheses to explain the observed data.
\begin{enumerate}[label=\alph*)]
  \item Suppose that the data is generated by a causal model with observable graph depicted in Figure~\ref{fig:simpson_paradox_causal_hypotheses}(a). Derive an expression for
    $$P(Y=1 \given \doit(X=x))$$
    in terms of the observed distribution $P(X,Y,Z)$, if possible.
    If you believe that such an expression does not exist, it suffices to state this without proof.
    Do the same for
    $$P(Y=1 \given Z=z,\doit(X=x)).$$
    Would you advise a patient to take the drug? Would your advice depend on the sex of the patient? 
    Motivate your answer numerically using the expressions just derived.
    \score{5}
    \answer{This can be ansered using the back-door criterion, or using the do-calculus.
    The back-door criterion immediately yields an expression, noting that $\{Z\}$ is the valid adjustment set.
    $$P(Y=1 \given \doit(X=x)) = \sum_{z \in \C{Z}} P(Y=1 \given Z=z, X=x) P(Z=z) \text{ a.s.}$$
    This can also be obtained from applying the chain rule and the do-calculus with $Y \Perp I_X \given X, Z$ and $Z \Perp I_X$.
    The conditional back-door criterion (with empty adjustment set), or the do-calculus with $Y \Perp I_X \given X, Z$, gives:
    $$P(Y=1 \given Z=z,\doit(X=x)) = P(Y=1 \given Z=z,X=x) \text{ a.s.}$$
    These conditional probabilities correspond with the recovery rates in the sex-specific tables. 
    Hence, patients of both sexes are better off by \emph{not} taking the drug.
    This implies that also a patient of unknown sex is better off by \emph{not} taking the drug.}
    \points{2 pt for the derivation of the causal effect, 2 pt for the derivation of the conditional causal effect, 1 pt for a well-motivated sensible advice.}
  \item The same question for hypothesis Figure~\ref{fig:simpson_paradox_causal_hypotheses}(b).
    \score{1}
    \answer{The same answer as for (a).}
  \item The same question for hypothesis Figure~\ref{fig:simpson_paradox_causal_hypotheses}(c).
    \score{5}
    \answer{Under this hypothesis, $\{Z\}$ is no longer a valid adjustment set, but $\emptyset$ is:
    $$P(Y=1 \given \doit(X=x)) = P(Y=1 \given X=x) \text{ a.s.}$$
    This also follows from the do-calculus with $Y \Perp I_X \given X$.
    These conditional probabilities correspond with the recovery rates in the aggregated table.
    Hence, for a patient of unknown sex the rational choice is to take the drug.
    There exists no valid adjustment set for the conditional treatment effect $P(Y=1 \given Z=z,\doit(X=x))$.
    Also the do-calculus will not help to identify this quantity (it is not identifiable from the data).
    Therefore, it seems that sex-specific advice would be hard to give based on these data
    (Perhaps one could derive lower and upper bounds on these probabilities, though).}
  \item The same question for hypothesis Figure~\ref{fig:simpson_paradox_causal_hypotheses}(d).
    \score{5}
    \answer{The empty set is still valid for adjustment:
    $$P(Y=1 \given \doit(X=x)) = P(Y=1 \given X=x) \text{ a.s.}$$
    which can also be derived from $Y \Perp I_X \given X$.
    Hence, for a patient of unknown sex the rational choice is to take the drug.
    In this case, we can also estimate the conditional treatment effect:
    $$P(Y=1 \given Z=z,\doit(X=x)) = P(Y=1 \given Z=z,X=x) \text{ a.s.}$$
    However, this may be misleading, because $Z$ is not a ``pre-treatment'' variable (but is caused by treatment under this hypothesis). 
    $Z$ now corresponds to sex \emph{after} treatment, which may differ from sex before treatment
    (perhaps the drug has the side effect of changing sex in some patients). 
    Na\"ively interpreting this expression counterfactually seems to tell us that for those patients that are male or female after the treatment, they would have been better off without he drug.
    Yet, we would advice a patient of unknown sex to take the drug.
    That doesn't make sense!
    So what looks mathematically like a `conditional treatment effect' does not tell anything useful for deciding on optimal actions under this hypothesis: it would compare apples with oranges.
    If we wish to give sex-specific advice on whether someone should take the drug, we should condition on sex \emph{before} treatment, which is not available in the data (at least under this causal hypothesis, where $Z$ is an effect of $X$).
    So what we can conclude from the data under this hypothesis is only that recovery will be higher if we treat the entire population than if we treat noone.}
  \item The same question for hypothesis Figure~\ref{fig:simpson_paradox_causal_hypotheses}(e).
    \score{2}
    \answer{Under this hypothesis, both quantities cannot be identified (at best, they can be bounded).
    Since the bounds will be non-informative (really?), no advice can be given here.}
  \item Would you consider all hypotheses (a)-(e) equally likely, or would you reject some of these \emph{a priori} based on common sense background knowledge about causal relations?
  \score{2}
  \answer{A reasonable assumption seems to be that treatment $X$ does not cause sex $Z$, allowing to reject (d).
  If we rule out selection bias in the past (the male/female ratio in the data is 50\%, so the data do not contain evidence that there has been selection bias), then it also seems unlikely that there are unobserved common causes of $Z$ and $X$, and of $Z$ and $Y$.
  This would rule out hypotheses (b) and (c). 
  But the validity of that assumption is debatable (and untestable from the available data).
  }
  \points{If a student interprets a bidirected edge $A \huh B$ as that $A$ does not cause $B$, zero points (that is so wrong). TODO: rephrase the question (a student answered just: ``I wouldn't consider them equally likeley'').}
  \item A law suit is filed against the pharmaceutical company that developed the drug, because several male patients who took the drug died quickly afterwards.
    The plaintiff claims that she inferred from the data that the probability that these patients would have recovered (and hence still be alive) had they not been treated, is 70\%.
    The court asks you for your expert advice in the matter. 
    What would be your response?
  \score{2}
  \answer{Such counterfactual quantities are generically unidentifiable from observed data, without making very strong modeling assumptions.
    As reliable models probably don't exist (given that the disease is new and the drug is experimental), the inferred probability has to be taken with a large grain of salt.
    At best, the data may lead to a lower and an upper bound on the counterfactual probability in the claim, instead of a point estimate as provided by the data scientist.}
    \points{The answer ``the claim is not founded on a causal model'' is worth 1 point.}
  \item Assume that treatment $X$ does not cause sex $Z$. 
    \score{3}
    Prove that Simpson's paradox cannot occur if the data come from a proper randomized controlled trial, in which treatment $X$ is completely randomized.
    \answer{We model $X$ as an exogenous input variable.
    We may assume the following CADMG (where $Z \huh Y$ and $Z \hut Y$ are unlikely, but leaving them in the model doesn't harm).
    \begin{center}\begin{tikzpicture}
    \begin{scope}[xshift=8cm,yshift=-3cm]
  %      \draw[shadow] (-1,-2) rectangle (3,2);
      \node[ndint] (X) at (0,0) {$X$};
      \node[ndout] (Y) at (2,0) {$Y$};
      \node[ndout] (Z) at (1,1) {$Z$};
  %    \node[ndlat] (L) at (1,-1) {$L$};
      \draw[tuh] (X) to (Y);
  %    \draw[tuh] (Z) to (X);
      \draw[tuh] (Z) to (Y);
  %    \draw[tuh] (L) to (X);
  %    \draw[tuh] (L) to (Y);
      \draw[huh,bend right] (Z) edge (Y);
      \draw[tuh,bend right] (Y) edge (Z);
    \end{scope}\end{tikzpicture}\end{center}
%Suppose treatment is beneficial for both sexes.
%$$P(Y_1 = 1 \given Z=z) > P(Y_0 = 1 \given Z=z)$$
%for all $z$ implies
%$$P(Y_1 = 1) = \sum_{z} P(Y_1 = 1 \given Z=z) P(Z=z)
%             > \sum_{z} P(Y_0 = 1 \given Z=z) P(Z=z)
%             = P(Y_0 = 1)$$
%Hence, treatment is beneficial overall.
%
Now
$$P(Y = 1 \given Z=z, \doit(X=1)) > P(Y = 1 \given Z=z, \doit(X=0)) \text{ a.s.}$$
for all $z$ implies
  \begin{equation*}\begin{split}
  P(Y = 1 \given \doit(X=1)) 
    &= \sum_{z} P(Y = 1 \given Z=z, \doit(X=1)) P(Z=z \given \doit(X=1)) \\
    &= \sum_{z} P(Y = 1 \given Z=z, \doit(X=1)) P(Z=z \given \doit(X=0)) \\
    &> \sum_{z} P(Y = 1 \given Z=z, \doit(X=0)) P(Z=z \given \doit(X=0)) \\
    &= P(Y = 1 \given \doit(X=0)) \text{ a.s.}
  \end{split}\end{equation*}
Here we first used the chain rule, and then $Z \Indep X$.
  The latter follows from the assumptions ($X$ is completely randomized and does not cause $Z$).
  Hence, the weird reversal of which treatment is best under additional conditioning on sex cannot occur.
%
%On the other hand,
%$$P(Y = 1 \given X=1, Z=z) > P(Y = 1 \given X=0, Z=z)$$
%for all $z$ does \emph{not} imply
%  \begin{equation*}\begin{split}
%  P(Y = 1 \given X=1) 
%    &= \sum_{z} P(Y = 1 \given X=1, Z=z) P(Z=z \given X=1) \\
%    &\ne \sum_{z} P(Y = 1 \given X=1, Z=z) P(Z=z \given X=0) \\
%    &> \sum_{z} P(Y = 1 \given X=0, Z=z) P(Z=z \given X=0) \\
%    &= P(Y = 1 \given X=1)
%  \end{split}\end{equation*}
  }
\end{enumerate}
